\documentclass[twocolumn,10pt]{article}

\usepackage[margin=0.75in]{geometry}
\usepackage{amsmath,amssymb}
\usepackage{graphicx}
\usepackage{xcolor}
\usepackage[colorlinks=true,linkcolor=blue,citecolor=blue,urlcolor=blue]{hyperref}
\usepackage{authblk}
\usepackage{abstract}

\renewcommand{\Affilfont}{\small}
\setlength{\affilsep}{0.3em}

% provisional-number flag
\newcommand{\prov}[1]{\textcolor{red}{#1}}

\title{\vspace{-2em}\textbf{Which generalized entropy? Discrimination among horizon-entropy
dark-energy models with DESI DR2 and Planck, and the closure of the
structure-growth channel}}

\author[]{Eric Schultz}
\affil[]{Independent researcher \quad ORCID: 0009-0006-6283-1696}
\date{\today}

\begin{document}

\twocolumn[
  \begin{@twocolumnfalse}
  \maketitle
  \begin{abstract}
  \noindent
  Horizon-entropy (holographic) dark energy is a large and active post-DESI
  program, but the literature is fragmented: each study fits a single entropy
  prescription and reports that it can reproduce the data, leaving open whether
  the data \emph{discriminate} between prescriptions. Using a uniform pipeline
  applied to real DESI DR2 baryon-acoustic-oscillation data and Planck cosmic
  microwave background distance priors, we show that they do. Casting Barrow and
  Tsallis holographic dark energy as a single power-law family
  $\rho_{\rm DE}\propto L^{n}$ and treating R\'enyi holographic dark energy as its
  genuinely distinct coupled generalization, we find that the \emph{direction} of
  entropy deformation is decisive: steeper $\rho(L)$ (Tsallis $\delta_T<1$)
  relieves the tension of standard holographic dark energy with the data, while
  flatter $\rho(L)$ (Barrow $\Delta>0$) worsens it; no variant is preferred over
  $\Lambda$CDM. We further show, analytically, that the linear structure-growth
  rate carries no information about these models beyond the background expansion:
  for any event-horizon holographic dark energy, sub-horizon dark-energy
  perturbations are suppressed relative to matter as $1/k^{3}$ by phase-averaging
  of the gravitational potential along the future light cone, so $f\sigma_8$
  cannot separate the models from one another or from $\Lambda$CDM beyond what
  background geometry already achieves. The proposed structure-growth
  confrontation of these models is therefore, for this class, structurally empty.
  \vspace{1em}
  \end{abstract}
  \end{@twocolumnfalse}
]

\section{Introduction}
\label{sec:intro}

The thermodynamic route to cosmology---in which the Friedmann equation is
recovered as the first law of thermodynamics applied to the apparent
horizon---has generated a broad family of dark-energy models built from
generalized (non-additive) entropies. Replacing the Bekenstein--Hawking
area law $S\propto A$ with a deformed law (Barrow, Tsallis, R\'enyi,
Kaniadakis, Sharma--Mittal, \dots) and identifying an infrared cutoff $L$
yields a holographic dark-energy (HDE) density $\rho_{\rm DE}\sim L^{n}$ or a
generalization thereof~\cite{Li2004,Wang2017,Barrow2020,Saridakis2020}. Post-DESI,
this has become a crowded program.

The crowding, however, hides a gap. Individual studies typically fix one
entropy prescription, fit it, and report that it can accommodate the data.
What is missing is a \emph{uniform} treatment that asks the comparative
question: given the same data and the same pipeline, do the data prefer one
entropy deformation over another, or are they all equivalent phenomenological
fits? We answer this here.

A second, sharper motivation comes from the fitting literature itself.
Confronting Kaniadakis holographic dark energy with DESI DR2, cosmic
chronometers, and supernovae, Luciano et al.~\cite{Luciano2025,Luciano2026}
find a structural degeneracy between the holographic parameters and the matter
density at the background level, and conclude explicitly that expansion data
alone are insufficient to discriminate the entropy deformation, so that
\emph{perturbation observables are essential}. This is precisely the call our
second result addresses. We show that, for event-horizon holographic dark
energy, the structure-growth channel carries no information beyond the
background---so the perturbation observables the field is turning to are, for
this class, empty of discriminating power.

This paper is a companion to Ref.~\cite{StepPaper}, which developed a
detectability and exclusion framework for localized features in the dark-energy
expansion history and, as a worked example on a physical theory, constrained
the Barrow deformation of standard HDE. We inherit that paper's validated
CMB distance-prior likelihood and its framing of what expansion-history
geometry can and cannot resolve, and we extend the analysis from a single
model to the model \emph{family}, adding two results. First, a background-level
discrimination: the data pick a side (Sec.~\ref{sec:discrim}). Second, and
independently, a negative result at the perturbation level: the structure-growth
channel is closed for the whole family (Sec.~\ref{sec:pert}). Together these
complete the picture---the models are fully testable with present data at the
background level, and growth data add nothing.

Throughout we adopt a spatially flat universe with
$\Omega_{m0}=0.31$ unless stated otherwise, and we use natural units in which
$3/(8\pi G)=H_0=1$ where convenient.

%========================================================================
\section{A unified framework for horizon-entropy dark energy}
\label{sec:framework}

\subsection{Barrow and Tsallis as one power law}

Barrow entropy, $S\propto A^{1+\Delta/2}$ with $\Delta\in[0,1]$ a
quantum-gravitational deformation parameter, and Tsallis entropy,
$S\propto A^{\delta_T}$, both lead---through the holographic
construction with a future-event-horizon cutoff---to a power-law dark-energy
density
\begin{equation}
\rho_{\rm DE}\;\propto\; L^{\,n}, \qquad L = R_h = a\!\int_t^\infty\!\frac{dt'}{a(t')},
\label{eq:powerlaw}
\end{equation}
with exponent mappings
\begin{equation}
n_{\rm Barrow}=\Delta-2, \qquad n_{\rm Tsallis}=2\delta_T-4 .
\label{eq:mappings}
\end{equation}
Both reduce to standard HDE~\cite{Li2004} at $n=-2$ (Barrow $\Delta=0$,
Tsallis $\delta_T=1$). The two ``different'' entropies are therefore the same
one-parameter family, seen from opposite sides of the standard-HDE point:
Barrow explores $n>-2$ (flatter $\rho(L)$), Tsallis $\delta_T<1$ explores
$n<-2$ (steeper $\rho(L)$).

We stress that this equivalence is asserted at the level of the resulting
dark-energy \emph{density scaling} $\rho_{\rm DE}\propto L^n$, which is what the
background expansion and all our observables depend on---not necessarily at the
level of the horizon first law. Barrow and Tsallis entropies are distinct
functionals, and a full thermodynamic derivation from each (through the
apparent-horizon first law $dE=TdS$) need not yield identical modified
Friedmann equations term by term. What we use, and all that the fits require,
is that with a future-event-horizon cutoff both produce a power-law density
with the exponent mappings above; two models with the same $n$ share the same
background dynamics and are observationally identical here regardless of the
entropy functional they descend from. The unification is thus phenomenological
and exact at the level of the observable, and we make no claim of equivalence
at the level of the underlying field theory.

A single background equation covers both. Writing $x=\ln a$ and
$\Omega\equiv\Omega_{\rm DE}$,
\begin{equation}
\frac{d\Omega}{dx}=\Omega(1-\Omega)\left[1-\frac{n}{c}\,\Omega^{-1/n}\right],
\label{eq:unified-ode}
\end{equation}
with $c$ the standard HDE constant. Equation~\eqref{eq:unified-ode} reproduces
the original Barrow equation of motion to machine precision, and at $n=-2$,
$c=1$, $\Omega_{\rm DE0}=0.69$ yields the exact standard-HDE value
$w_0=-0.891$. The corresponding equation of state follows analytically from
Eq.~\eqref{eq:unified-ode} and continuity,
\begin{equation}
w(\Omega)=\frac{-1+(n/c)\,\Omega^{-1/n}}{3},
\label{eq:eos}
\end{equation}
with the correct limits $w\to-1/3$ as $\Omega\to0$ (matter era) and $w\to-1$
as $\Omega\to1$ (de Sitter future).

A numerical remark, relevant to reproducibility: the future event horizon
makes de Sitter a \emph{forward} attractor, so integrating
Eq.~\eqref{eq:unified-ode} forward from today is unstable (errors grow toward
$\Omega\to1$). We integrate instead by shooting from a modest future point
$x\simeq3$ and solving both directions, tuning $\Omega(x_{\rm start})$ to hit
$\Omega_{\rm DE0}$; and we obtain $w$ from the analytic form
Eq.~\eqref{eq:eos} rather than by numerical differentiation.

\subsection{R\'enyi as a genuine generalization}

R\'enyi holographic dark energy does not reduce to a power law. With
\begin{equation}
\rho_{\rm DE}=\frac{3C^2}{8\pi L^2\,(1+\delta\pi L^2)},
\label{eq:renyi}
\end{equation}
the cutoff $L$ does not drop out of the background dynamics, and one is left
with a genuinely coupled two-dimensional system in $(\Omega_{\rm DE},L)$. We
solve it by backward integration from the de Sitter future attractor,
shooting on the normalization to hit $\Omega_{\rm DE0}$; in the
$\delta\to0$ limit this converges to standard HDE ($w_0=-0.889$ at
$\delta=0.01$). Here $C^2$ plays the role of the squared HDE constant,
$C^2=c^2$, so that $C^2=1$ recovers $c=1$ and $w_0\simeq-0.891$; we adopt this
normalization to match the power-law family.

Although R\'enyi is not a power law, its perturbations reduce to an
\emph{effective local exponent}, a fact we exploit in
Sec.~\ref{sec:pert}: writing $\rho_{\rm DE}=f(L)$,
\begin{equation}
n_{\rm eff}(L)\equiv\frac{L f'(L)}{f(L)}=-2-\frac{2\delta\pi L^2}{1+\delta\pi L^2},
\label{eq:neff}
\end{equation}
which runs monotonically from $-2$ (small $L$, standard HDE) to $-4$
(large $L$, R\'enyi-dominated) and is bounded in $[-4,-2]$.

\subsection{Kaniadakis as a fourth member}
\label{sec:kaniadakis}

Kaniadakis holographic dark energy, built from the relativistic
$\kappa$-deformed entropy~\cite{Luciano2025}, is a further one-parameter
extension. Like R\'enyi it is not a power law: the $\kappa$-deformation of the
area law modifies the Friedmann equation and yields a coupled background whose
equation of state can be quintessence-like, phantom-like, or cross the phantom
divide, with the acceleration transition at $z\simeq0.6$. It has been shown to
be a competitive but not preferred alternative to $\Lambda$CDM against DESI
DR2, cosmic chronometers, and supernovae~\cite{Luciano2025,Luciano2026}. We
include it as a fourth member of the family; being a distinct coupled
background it is treated like R\'enyi rather than through the power-law ODE
Eq.~\eqref{eq:unified-ode}. Crucially for Sec.~\ref{sec:pert}, Kaniadakis
holographic dark energy uses the future event horizon as its infrared cutoff,
so---whatever its background behavior---its perturbations fall under the same
closed-channel theorem: the entropy deformation changes the density law, not
the null-geodesic definition of the horizon.

%========================================================================
\section{The data discriminate at the background level}
\label{sec:discrim}

\subsection{Data and likelihood}

We fit each model to real DESI DR2 anisotropic BAO
(six tracers over $z=0.51$--$2.33$, $D_M/r_d$ and $D_H/r_d$ with their
correlations)~\cite{DESIDR2} and Planck 2018 CMB distance
priors~\cite{Chen2019,Planck2018}, the compressed likelihood in
$\{R,\ell_A,\omega_b\}$ with the Chen et al.\ inverse covariance. The shift
parameter, acoustic scale, and sound horizon are
\begin{align}
R &= \sqrt{\Omega_{m}}\,\textstyle\int_0^{z_*}\!\frac{dz}{E(z)}, \qquad
\ell_A = \pi\,\frac{D_M(z_*)}{r_s(z_*)},\\
r_s(z_*) &= \textstyle\int_{z_*}^{\infty}\!\frac{c_s(z)}{E(z)}\,dz, \quad
c_s=\frac{1}{\sqrt{3(1+R_b)}},
\end{align}
with $R_b=3\omega_b/(4\omega_\gamma)/(1+z)$ and $z_*$ from the
Hu--Sugiyama fitting formula. We compute $r_s$ using the analytic
matter-plus-radiation $E(z)$ above $z_*$, where dark energy is dynamically
negligible; this is essential, as retaining the residual dark-energy term
makes the high-redshift $r_s$ integral converge incorrectly. The likelihood
is validated in Ref.~\cite{StepPaper}: pure-$\Lambda$CDM Planck parameters
reproduce $R=1.750$ and $\ell_A=300.4$ (to $0.3\%$).

We profile over $(\Omega_m,H_0,c)$ and, for the final results, over
$\omega_b$ as well (Sec.~\ref{sec:disc-caveat}).

\subsection{R\'enyi at the CMB}

Extending R\'enyi HDE to $z_*\simeq1090$ requires care. The future-event-horizon
variable $L\to0$ at high redshift, so the R\'enyi density
Eq.~\eqref{eq:renyi} grows in absolute terms ($\rho_{\rm DE}\sim C^2/L^2\sim
(1+z)^2$). This is not a pathology: since matter scales as $(1+z)^3$ and
radiation as $(1+z)^4$, the fractional contribution $\Omega_{\rm DE}(z)\to0$
monotonically ($\Omega_{\rm DE}\sim2\times10^{-4}$ at $z_*$), so the model has a
small, integrable early dark-energy tail rather than a divergence in any
observable. What the growing $\rho_{\rm DE}$ does cause is amplified numerical
noise if the horizon system is integrated naively through the near-singular
regime. We handle it with a matching redshift $z_{\rm match}$: below it we
evolve the full coupled R\'enyi system; above it we use the analytic
matter-plus-radiation $E(z)$ \emph{retaining} the $(1+z)^2$ dark-energy tail
rather than setting it to zero, so no physical early-DE contribution to
$r_s(z_*)$ is neglected. The split is therefore a numerical device, not a
physical transition. The result is insensitive to $z_{\rm match}$ once the tail
is carried; dropping it instead (an earlier, cruder choice) leaves a
$z_{\rm match}$-dependent residual of $\sim0.2$ in $\ell_A$ that converges away
for $z_{\rm match}\gtrsim50$. The procedure reproduces $z_*=1091.8$ and, in the
$\delta\to0$ limit with $C^2=1$, recovers standard HDE ($w_0=-0.887$,
$R=1.70$, $\ell_A=300.1$). The full numerical mechanics are given in the
Appendix.

\subsection{Discrimination result}
\label{sec:discrim-result}

Table~\ref{tab:chi2} collects the best-fit total $\chi^2$ (DESI BAO $+$ CMB
distance priors, $H_0$ free) across the family. Three statements stand out,
none of which is available from single-model studies.

\emph{(i) Standard HDE is already disfavored, before any deformation.} With a
future-event-horizon cutoff, standard HDE ($n=-2$) sits in strong tension with
DESI DR2 BAO geometry: $\chi^2_{\rm BAO}\simeq\prov{47}$ versus
$\simeq\prov{10}$ for $\Lambda$CDM, a $\Delta\chi^2\simeq\prov{37}$ penalty
independent of the CMB. This is consistent with the post-DESI finding that
standard HDE is strained.

\emph{(ii) The direction of deformation is decisive.} Steepening $\rho(L)$
(Tsallis $\delta_T<1$, $n<-2$) \emph{relieves} the tension, primarily by
repairing the CMB fit; flattening it (Barrow $\Delta>0$, $n>-2$) makes matters
worse. The best HDE variant is Tsallis $\delta_T\simeq\prov{0.64}$, with total
$\chi^2\simeq\prov{24.5}$ (from $\prov{29}\to\prov{2.6}$ in the CMB), while
Barrow $\Delta=0.1$ rises to $\simeq\prov{43.8}$.

\emph{(iii) None beats $\Lambda$CDM.} The best HDE fit remains
$\Delta\chi^2\simeq\prov{14}$ above $\Lambda$CDM. Generalized entropy is thus
not monolithic: the data select \emph{steeper} $\rho(L)$ but do not favor any
member of the family over the cosmological constant.

\begin{table}[t]
\caption{Best-fit total $\chi^2$ (DESI DR2 BAO $+$ Planck CMB distance priors,
$H_0$ free) across the horizon-entropy family. \textbf{Numbers in
\prov{red} are provisional}, carried over from an earlier fixed-$\omega_b$ run,
and are to be regenerated with the $\omega_b$-marginalized pipeline
(Sec.~\ref{sec:disc-caveat}) before submission. Qualitative ordering
is expected to be robust; absolute values are not yet final.}
\label{tab:chi2}
\begin{center}\small
\begin{tabular}{lcc}\hline\hline
Model & deformation & total $\chi^2$ \\
\hline
\hline
$\Lambda$CDM & --- & \prov{$\sim$10} \\
Tsallis & $\delta_T=0.64$ & \prov{24.5} \\
Tsallis & $\delta_T=0.70$ & \prov{27.3} \\
Standard HDE & $n=-2$ & \prov{38.1} \\
Barrow & $\Delta=0.1$ & \prov{43.8} \\
R\'enyi & $\delta=0.1$ & \prov{BAO $\sim$11; CMB via split} \\
Kaniadakis & $K$ small & \prov{competitive; $\Lambda$CDM favored} \\
\hline\hline
\end{tabular}
\end{center}
\end{table}

\subsection{The $H_0$ preference}

Freeing $H_0$ improves every HDE fit substantially, toward the local
distance-ladder value $H_0\simeq74$: standard HDE moves from
$\chi^2\simeq\prov{75.5}$ at $H_0=67.4$ to $\simeq\prov{38.1}$ at
$H_0\simeq\prov{74}$. This is a known property of future-event-horizon HDE and
is Hubble-tension-adjacent. We flag one honest caveat: our CMB enters only as
distance priors, which are more permissive in $H_0$ than the full Planck
likelihood, so the strength of the $H_0\simeq74$ preference is partly a
property of the compressed likelihood. The \emph{relative} model ordering is
robust to this.

\subsection{Status of the numbers}
\label{sec:disc-caveat}

The $\chi^2$ values in Table~\ref{tab:chi2} are provisional. They derive from a
run in which $\omega_b$ was fixed at $0.02237$. Because $\sigma(\ell_A)$ is
small, the CMB likelihood is strongly sensitive to $\omega_b$ through $r_s$: a
model preferring $\omega_b$ shifted by as little as $\sim4\times10^{-4}$---well
within the Planck uncertainty---incurs a spurious penalty of tens in $\chi^2$
when $\omega_b$ is held fixed. This penalty is \emph{not} uniform across models,
since different deformations reshape the low-redshift distance and therefore
pull $\omega_b$ in different directions. The final table must therefore profile
(marginalize) over $\omega_b$ per model. We expect the qualitative ordering
(steeper helps, flatter hurts, none beats $\Lambda$CDM) to survive, but the
absolute values, and in principle the ordering, must be confirmed with the
$\omega_b$-marginalized pipeline. The $\omega_b$ profiling itself is a
one-dimensional minimization inside each model's $\chi^2$, validated on
self-consistent mocks; what remains is to run it on the fully calibrated
distance-prior pipeline of Ref.~\cite{StepPaper}.

%========================================================================
\section{The structure-growth channel is closed}
\label{sec:pert}

Recent analyses call for confronting horizon-entropy dark energy at the
perturbative, structure-growth level, on the expectation that $f\sigma_8$
adds information beyond the background. We show analytically that for
event-horizon HDE it does not: the models are, to the precision that matters,
smooth, and the growth rate is degenerate with $H(z)$.

\subsection{Growth depends on the cosmology only through $H(a)$}

For smooth (non-clustering) dark energy the linear growth factor obeys, in
$e$-folds $N=\ln a$,
\begin{equation}
D''+\left(2+\frac{H'}{H}\right)D'-\frac{3}{2}\,\Omega_m(a)\,D=0 .
\label{eq:growth}
\end{equation}
Every cosmology-dependent term is a function of $H(a)$ alone: both the
friction $2+H'/H$ and the source $\tfrac32\Omega_m(a)=\tfrac32
\Omega_{m0}a^{-3}/(H/H_0)^2$. There is no independent dark-energy degree of
freedom. Two models with identical $H(a)$ produce identical $D(a)$, hence
identical $f\sigma_8(z)$. Numerically, a CPL model tuned to mimic $\Lambda$CDM
to within $0.5\%$ in $H(z)$ over $0<z<2$ differs by at most $0.5\%$ in
$f\sigma_8$---the same order as the background difference, far below current
$\sim5$--$10\%$ growth-rate errors. The growth observable inherits the
background (non-)difference and adds nothing.

\subsection{The sound speed cannot be adiabatic}

Independent information would require dark energy to cluster, i.e.\ a rest-frame
sound speed $c_s^2<1$. But the adiabatic value
$c_a^2=w-\dot w/[3H(1+w)]$ is negative or singular for $w<-1/3$---generic for
dark energy, and singular near the phantom line several variants approach---so
it must be discarded on stability grounds (imposing it drives sub-horizon
$\delta_{\rm DE}$ to blow up). HDE has no Lagrangian kinetic term from which to
read a rest-frame $c_s^2$; the density is defined non-locally by a horizon.
The effective sound speed is thus whatever the horizon prescription implies,
which we now compute.

\subsection{Horizon perturbations phase-average away}

Physically, because the future event horizon is defined by an integral along
the future light cone, its value at the observer sums the gravitational
potential over a long null path that threads many alternating crests and
troughs of a small-scale perturbation; these contributions largely cancel, so
the horizon---and the dark-energy density tied to it---barely responds to
sub-horizon structure. We now make this quantitative.

With $\rho_{\rm DE}\propto R_h^{\,n}$, exactly
$\delta_{\rm DE}=n\,(\delta R_h/R_h)$. The future event horizon is a null
construct; perturbing the radial null geodesic in Newtonian gauge, the comoving
horizon acquires
\begin{equation}
\delta\chi_k(t)=2\,e^{ikx}\!\int_0^{\chi_h}\!\Phi_k(u)\,e^{iku}\,du,
\qquad u=\chi(t,t') ,
\label{eq:dchi}
\end{equation}
the comoving distance along the future ray. Here $\Phi_k$ is the
Newtonian-gauge potential, which coincides with the Bardeen invariant
$\Phi_A$; since the null ray is itself a geometric (gauge-independent) object,
$\delta R_h/R_h$ is built entirely from gauge-invariant ingredients and is not
an artifact of the conformal Newtonian slicing. A residual gauge shift in the
dark-energy density contrast is $\mathcal{O}(aH/k)\,(1+w)\,v_{\rm DE}$; with
$v_{\rm DE}$ sourced at $\mathcal{O}(\Phi)$ this is of order $\delta_m/k^3$,
the same order as the result itself, so no slicing can promote $\delta_{\rm DE}$
to a non-decaying mode---there is none to transform into. For a sub-horizon
mode
($k\gg aH$) the kernel $e^{iku}$ oscillates across the integration range;
integration by parts with any smooth envelope gives
$\int_0^{\chi_h}\Phi_k\,e^{iku}\,du=\mathcal{O}(1/k)$. Hence
$|\delta_{\rm DE}|\sim(aH/k)|\Phi|$: the potential is rapidly
\emph{phase-averaged} along the future light cone, and a small-scale
fluctuation contributes essentially nothing to the non-local horizon integral.

Relative to the matter that actually clusters, the suppression is stronger
still. Sub-horizon, $\Phi$ is sourced by matter through Poisson,
$k^2\Phi=-4\pi G a^2\rho_m\delta_m$, so $\Phi\sim\delta_m/k^2$ and
\begin{equation}
\frac{\delta_{\rm DE}}{\delta_m}\;\sim\;\frac{1}{k^{3}}\qquad(k\gg aH).
\label{eq:k3}
\end{equation}
Dark energy is smooth to third order in $aH/k$ relative to matter. We record
this as the central result of this section:
\begin{quote}
\emph{No-clustering theorem for event-horizon dark energy.} For any
horizon-entropy dark-energy model with a future-event-horizon cutoff, the
sub-horizon dark-energy density contrast is suppressed as
$\delta_{\rm DE}/\delta_m\sim(aH/k)^3$; the linear growth of structure is
governed by $H(a)$ alone, and $f\sigma_8$ carries no information about the
model beyond its background expansion.
\end{quote}
(We state
the $1/k$ scaling of $\delta_{\rm DE}$ directly rather than dressing it as an
effective $c_s^2=1$; a genuine $c_s^2=1$ fluid gives a $1/k^2$ forced
amplitude, so the horizon suppression relative to $\Phi$ is in fact weaker than
$c_s^2=1$, while relative to matter it is $1/k^3$.)

The integration by parts deserves its boundary term made explicit, since it is
local and not phase-averaged. Writing
$\int_0^{\chi_h}\Phi_k e^{iku}du = [\Phi_k e^{iku}/ik]_0^{\chi_h}
- (1/ik)\int_0^{\chi_h}\Phi_k' e^{iku}du$, the surface piece at the observer,
$-\Phi_k(0)/ik$, does not oscillate in $k$ and so is not averaged away. But it
is itself $\mathcal{O}(1/k)$: $\Phi_k(0)$ is the finite local potential, not a
quantity that grows with $k$. The far endpoint $\Phi_k(\chi_h)e^{ik\chi_h}/ik$
oscillates in $k$ and averages down further. Both the surface and bulk
contributions are therefore $\mathcal{O}(1/k)$, so
$\delta R_h\sim\mathcal{O}(1/k)$ and the scaling of Eq.~\eqref{eq:k3} is
unchanged. A degradation to $1/k^2$ would require $\delta R_h=\mathcal{O}(1)$,
i.e. $\Phi_k(0)=\mathcal{O}(k)$, which is unphysical. A numerical evaluation of
Eq.~\eqref{eq:dchi} with a nonzero local potential $\Phi_k(0)\neq0$ confirms
$|\delta\chi_k|\propto k^{-1}$ exactly, the surface term supplying the leading
$1/k$ rather than any $1/k^0$ piece.

We keep the leading light-ray phase term; local metric corrections at the
observation point are also $\mathcal{O}(aH/k)$ and subleading. A three-closure
numerical integration across $k=1$--$1000$ confirms the picture: an imposed
$c_s^2=1$ leaves sub-horizon $\delta_{\rm DE}$ at $\sim10^{-9}$ (smooth), the
adiabatic $c_s^2<0$ blows up (unphysical), and the horizon closure tracks the
smooth branch.

\subsection{Extension to R\'enyi}

R\'enyi modifies the density law $f(L)$, not the definition of $L$: the cutoff
remains the future event horizon, a geometric null construct. Its perturbation
is $\delta_{\rm DE}=n_{\rm eff}(L)\,(\delta L/L)$ with $n_{\rm eff}$ the bounded
quantity of Eq.~\eqref{eq:neff}. The scale dependence lives entirely in
$\delta L$---the same oscillatory horizon integral, which is kinematic (it
arises from null propagation, $ds^2=0\Rightarrow dx=dt/a$) and untouched by the
coupled background. The coupled system alters only bounded, smooth,
non-oscillatory quantities (the horizon size, the potential envelope, and the
$n_{\rm eff}$ prefactor). Integration by parts still gives $\mathcal{O}(1/k)$,
so $|\delta_{\rm DE}|=|n_{\rm eff}|\,\mathcal{O}(aH/k)\to0$ sub-horizon.
R\'enyi is smooth on the same footing as the power-law family.

\subsection{Relation to recent sound-speed constraints}
\label{sec:pert-clustering}

It might appear that our conclusion is in tension with recent work that, for
the first time, constrains the clustering of dynamical dark energy from
current data. Combining DESI DR2 BAO with Planck 2018 and the Union3 supernovae,
Yang et al.~\cite{Yang2025} report a bound on the effective sound speed,
$\log_{10}c_s^2=-3.00^{+2.9}_{-0.99}$, in a $w_0w_a$CDM background within the
parameterized-post-Friedmann and effective-field-theory descriptions. Taken at
face value this says the data can now ``see'' dark-energy clustering; how can
the growth channel be closed?

There is no contradiction, and the distinction is the substance of our result.
That analysis treats $c_s^2$ as a \emph{free} parameter, made estimable by the
time-varying equation of state, which breaks the degeneracy between $(1+w)$ and
$c_s^2$. Its constraint is a statement about what the data could reveal
\emph{if} dark energy has an adjustable sound speed. Horizon-entropy dark
energy is not of that kind. Its density is fixed non-locally by a horizon, with
no independent stress sector to carry a free $c_s^2$; the effective sound speed
is not a parameter to be fit but is \emph{determined} by the horizon
prescription, and for an event-horizon cutoff that determination is the smooth
branch, Eq.~\eqref{eq:k3}. In short: the sound-speed constraint says
\emph{if it clusters, we can measure it}; our result says \emph{event-horizon
holographic dark energy does not cluster, so there is nothing to measure}. The
two are complementary halves of the same physical statement, and together they
sharpen it---the growth data that can now probe a free sound speed have no
purchase on a model whose sound speed the horizon has already fixed.

\subsection{A consistency check: the growth index}
\label{sec:pert-gamma}

The smooth limit makes a concrete, falsifiable prediction for the growth index
$\gamma$, defined by $f=\Omega_m^\gamma$. Mehrabi et al.~\cite{Mehrabi2015}
derive the asymptotic index for holographic dark energy in two distinct
regimes. For \emph{homogeneous} (non-clustering) holographic dark energy the
general result $\gamma_\infty=3(w_\infty-1)/(6w_\infty-5)$, evaluated at the
holographic asymptotic equation of state $w_\infty\to-1/3$, gives
\begin{equation}
\gamma_\infty^{\rm homog}\simeq\frac{4}{7},
\end{equation}
about $5\%$ above the $\Lambda$CDM value $6/11$. For \emph{fully clustered}
holographic dark energy a different source term applies and the index becomes
$\gamma_\infty^{\rm clust}\simeq3(1-c_{\rm eff}^2)/7$, which at $c_{\rm eff}^2=0$
gives $3/7$. These are two separate branches, not endpoints of one
interpolation.

Our result places event-horizon holographic dark energy squarely in the
homogeneous branch: the $1/k^3$ suppression means the dark-energy density
contrast never sources the matter growth equation ($\delta_{\rm DE}\simeq0$),
so growth is governed by $H(a)$ and $\Omega_m(a)$ alone. The asymptotic
homogeneous value $\gamma_\infty^{\rm homog}\simeq4/7$ is set by the high-$z$
holographic equation of state $w\to-1/3$; it is \emph{not} directly the
observable index, because in these models the late-time attractor is de Sitter
($w\to-1$), not $w=-1/3$. To obtain the observable prediction we integrate the
growth equation numerically along the event-horizon background
($w=-\tfrac13(1+2\sqrt{\Omega_{\rm DE}})$ for the $c=1$ cutoff), reading off
$\gamma(z)=\ln f/\ln\Omega_m$:
\begin{center}
\begin{tabular}{lccccc}
$z$ & $2.0$ & $1.0$ & $0.5$ & $0.3$ & $0.0$ \\ \hline
$\gamma(z)$ & $0.562$ & $0.560$ & $0.560$ & $0.560$ & $0.561$ \\
\end{tabular}
\end{center}
The observable index is remarkably flat at $\gamma\simeq0.56$ across the range
where growth is measured---close to but distinct from the asymptotic $4/7\simeq
0.571$, and sitting about $0.015$ above the $\Lambda$CDM value $6/11\simeq
0.545$. The prediction is thus model-external and falsifiable: the entire
event-horizon class predicts $\gamma\simeq0.56$, weakly $z$-dependent and
insensitive to the entropy choice (which affects only $H(a)$, already
constrained at the background level). A measurement pinning $\gamma$ to the
$\Lambda$CDM value $0.545$ or below would disfavour the class; a value near the
fully clustered $3/7\simeq0.43$ would indicate the perturbations do source
growth, contradicting the closed channel. Separating $0.56$ from $0.545$
requires the precision of a Euclid-era growth survey, but the boundary is
clean.

\subsection{The one condition}
\label{sec:pert-condition}

The derivation assumes the infrared cutoff is the \emph{future event horizon}.
If instead a variant uses the Hubble horizon $L=1/H$ (or a mixed cutoff), then
$\delta L=-\delta H/H^2$ is a \emph{local} quantity with no oscillatory kernel
and no $1/k$ suppression, and dark energy could genuinely cluster. The precise
statement is therefore: \emph{event-horizon horizon-entropy dark energy is
smooth on sub-horizon scales, and $f\sigma_8$ carries no information beyond
$H(z)$, if and only if the model uses a light-ray (event-horizon) cutoff.} Our
models use the future event horizon; the result holds. (Hubble-cutoff HDE
independently fails to produce late-time acceleration in the presence of dark
matter, so the event horizon is the standard choice.) A Hubble- or mixed-cutoff
construction is the one place in this analysis where the perturbation channel
could reopen---a well-defined open problem.

It is worth noting that the event-horizon cutoff on which our result depends is
also the one the data prefer. Combining expansion and growth-rate data across
the three most common cutoffs, Akhlaghi et al.~\cite{Akhlaghi2018} find the
Ricci and Granda--Oliveros prescriptions disfavored---they predict a
non-negligible early-time dark-energy density---while the future-event-horizon
cutoff remains consistent with the combined data. Our scope is therefore not a
convenience: the cutoff for which the growth channel is provably closed is the
same cutoff the observations single out.

%========================================================================
\section{Discussion}
\label{sec:discussion}

We have shown two things about horizon-entropy dark energy that a single-model
fit cannot reveal. At the background level, a uniform confrontation with real
DESI DR2 BAO and Planck distance priors \emph{discriminates}: steeper
$\rho(L)$ (Tsallis $\delta_T<1$) relieves the standard-HDE tension, flatter
$\rho(L)$ (Barrow $\Delta>0$) worsens it, and none of the family is preferred
over $\Lambda$CDM. At the perturbation level, the structure-growth channel is
\emph{closed}: for any event-horizon member, sub-horizon dark-energy
perturbations are suppressed as $1/k^3$ relative to matter by phase-averaging
along the future light cone, so $f\sigma_8$ is degenerate with the background
$H(z)$ and adds no discriminating power.

The two results are complementary. The first says the models are testable now,
at the background level, and shows what the test yields. The second says the
verdict cannot be overturned by adding growth data: one cannot rescue a
background-disfavored member by appealing to a distinctive clustering signature,
because there is none. For this class the proposed structure-growth
confrontation is structurally empty---a useful thing to know before investing
in it.

Two limitations bound these conclusions. The background numbers of
Sec.~\ref{sec:discrim} are provisional pending the $\omega_b$-marginalized
run (Sec.~\ref{sec:disc-caveat}); we expect the qualitative
ordering to survive but state plainly that it must be confirmed. And the
closed-channel result is conditional on an event-horizon cutoff
(Sec.~\ref{sec:pert-condition}); the Hubble-cutoff case is left open and is the
natural next target, being the one construction in which growth data could
become informative.

Two matters of scope should be stated explicitly. First, we treat these models
as \emph{late-time} modifications of the dark-energy sector: the entropy
deformation governs the infrared, horizon-scale dark energy, while
early-universe thermodynamics---recombination, the sound horizon $r_s(z_*)$,
and the thermal history---are taken to be standard. This is what licenses the
use of the Planck distance priors $(R,l_A,\omega_b)$, which are extracted
assuming standard pre-recombination physics. It is a genuine restriction: a
construction in which the modified entropy also deforms the apparent-horizon
thermodynamics of the radiation era would shift $r_s(z_*)$ and invalidate the
compressed likelihood, and such a model would require a full Boltzmann
treatment rather than distance priors. Our results apply to the
phenomenologically standard case where the deformation is confined to the dark
sector at low redshift; we do not claim them for early-modified variants.
Second, the phantom members (Tsallis $\delta_T<1$, $n<-2$) can cross $w=-1$;
where they approach a future Big Rip at finite $t_{\rm rip}$, the event-horizon
integral $L=a\int_t^{t_{\rm rip}}dt'/a$ terminates at $t_{\rm rip}$ rather than
at infinity, and the background shooting must impose the finite-time boundary
there. This is well-defined but changes the future boundary condition; we have
restricted the fitted range to members with a de Sitter future, and flag the
phantom/Big-Rip branch as requiring separate boundary treatment.

\section*{Acknowledgments}
This work uses public DESI DR2 BAO and Planck data products and builds on the
detectability framework and validated distance-prior likelihood of
Ref.~\cite{StepPaper}.

\appendix
\section{Numerical mechanics of the horizon integration}
\label{app:numerics}

We collect the implementation details here to keep the main text focused on
results. Two points matter for reproducibility.

\emph{Background shooting.} The future event horizon makes de Sitter a forward
attractor, so integrating Eq.~\eqref{eq:unified-ode} forward from today is
numerically unstable: errors grow toward $\Omega\to1$. We instead shoot from a
modest future point $x=\ln a\simeq3$, integrating in both directions and tuning
$\Omega(x_{\rm start})$ to hit the measured $\Omega_{\rm DE0}$ at $x=0$. The
equation of state is taken from the analytic form Eq.~\eqref{eq:eos} rather than
by numerically differentiating $\Omega(x)$, which avoids amplifying integration
noise.

\emph{R\'enyi high-redshift matching.} For the R\'enyi (and Kaniadakis) coupled
backgrounds the density $\rho_{\rm DE}\sim C^2/L^2\sim(1+z)^2$ grows toward high
redshift while $\Omega_{\rm DE}\to0$. Rather than integrate the stiff coupled
system through this regime, we match at $z_{\rm match}$: the full coupled system
below, and the analytic matter-plus-radiation $E(z)$ above, \emph{retaining} the
$(1+z)^2$ dark-energy tail so no physical early-DE contribution to $r_s(z_*)$ is
dropped. The sound-horizon integral uses matter and radiation only for the
acoustic sound speed, $c_s=[3(1+R_b)]^{-1/2}$ with
$R_b=3\omega_b/4\omega_\gamma/(1+z)$, and $z_*$ from the Hu--Sugiyama fitting
formula. With the tail retained the compressed observables are stable against
$z_{\rm match}$; we use $z_{\rm match}\gtrsim50$ and verify convergence.





\begin{thebibliography}{9}

\bibitem{Li2004} M. Li, ``A Model of Holographic Dark Energy,''
Phys. Lett. B \textbf{603}, 1 (2004), arXiv:hep-th/0403127.

\bibitem{Wang2017} S. Wang, Y. Wang, and M. Li, ``Holographic Dark Energy,''
Phys. Rept. \textbf{696}, 1 (2017), arXiv:1612.00345.

\bibitem{Barrow2020} J. D. Barrow, ``The Area of a Rough Black Hole,''
Phys. Lett. B \textbf{808}, 135643 (2020), arXiv:2004.09444.

\bibitem{Saridakis2020} E. N. Saridakis, ``Barrow holographic dark energy,''
Phys. Rev. D \textbf{102}, 123525 (2020), arXiv:2005.04115.

\bibitem{Chen2019} L. Chen, Q.-G. Huang, and K. Wang,
``Distance Priors from Planck Final Release,''
JCAP \textbf{02}, 028 (2019), arXiv:1808.05724.

\bibitem{Planck2018} Planck Collaboration (N. Aghanim et al.),
``Planck 2018 results. VI. Cosmological parameters,''
Astron. Astrophys. \textbf{641}, A6 (2020)
[Erratum: Astron. Astrophys. \textbf{652}, C4 (2021)], arXiv:1807.06209.

\bibitem{DESIDR2} DESI Collaboration (M. Abdul Karim et al.),
``DESI DR2 Results II: Measurements of Baryon Acoustic Oscillations
and Cosmological Constraints,'' Phys. Rev. D \textbf{112}, 083515 (2025),
arXiv:2503.14738.

\bibitem{Luciano2025} G. G. Luciano et al., ``Late-Time Cosmological
Constraints on Kaniadakis Holographic Dark Energy,''
Eur. Phys. J. C \textbf{85}, 1384 (2025), arXiv:2509.17527.

\bibitem{Luciano2026} G. G. Luciano et al., ``Infrared Corrections and
Horizon Phase Transitions in Kaniadakis-Based Holographic Dark Energy,''
arXiv:2603.21218 (2026).

\bibitem{Yang2025} Y. Yang et al., ``Probing the sound speed and clustering
of dark energy: observational constraints from DESI, CMB, and supernovae,''
arXiv:2511.22478 (2025).

\bibitem{Mehrabi2015} A. Mehrabi, S. Basilakos, M. Malekjani, and Z. Davari,
``Growth of matter perturbations in clustered holographic dark energy
cosmologies,'' Phys. Rev. D \textbf{92}, 123513 (2015), arXiv:1510.03996.

\bibitem{Akhlaghi2018} I. A. Akhlaghi, M. Malekjani, S. Basilakos, and H. Haghi,
``Model selection and constraints from holographic dark energy scenarios,''
Mon. Not. R. Astron. Soc. \textbf{477}, 3659 (2018), arXiv:1804.02989.

\bibitem{StepPaper} E. Schultz, ``Detectability and exclusion of localized
dark-energy features in the expansion history,'' companion paper,
in preparation (2026).

\end{thebibliography}

\end{document}
